% SPDX-License-Identifier: 0BSD
\section{\texorpdfstring{Shared-Current DAG and Numerical Execution}{Shared-Current DAG and Numerical Execution}}
\label{sec:zgg-dag}

The current recursion becomes a directed acyclic graph (DAG) when physically
identical intermediate currents are shared. This section summarizes the
identity and ordering needed for that reuse and relates them to the three
execution modes.

\subsection{Current identity and reuse}

Two contributions share a current only when every physics-relevant label
agrees. Schematically,
\begin{equation}
 \Omega(J)=
 \bigl(t,I,H,s,F,Q,C,O,a,\pi_I\bigr),
 \label{eq:zgg-current-key}
\end{equation}
where the entries identify the particle species, external subset, helicity and
spin ancestry, flavour and additive charges, colour state, perturbative order,
auxiliary-current type, and colour-sensitive ordering. The subset determines
the current momentum \(q_I=\sum_{i\in I}q_i\), while the ordered labels prevent
distinct colour traversals from being merged.

Exact identity permits one stored current to serve every amplitude contribution
that needs it. A second kind of reuse applies when different currents require
the same numerical kernel value up to a known exact factor. Neither form of
reuse merges distinct physical helicity or colour-flow components.

\subsection{Topological ordering}

The graph is evaluated in increasing external-subset size. Every input is
therefore available before the current that consumes it. For the process in
Sec.~\ref{sec:zgg-example}, the ordering is schematically
\begin{equation}
 \boxed{\text{sources}}
 \longrightarrow
 \boxed{|I|=2}
 \longrightarrow
 \boxed{|I|=3}
 \longrightarrow
 \boxed{|I|=4}
 \longrightarrow
 \boxed{\text{amplitude closures}} .
 \label{eq:zgg-stage-chain}
\end{equation}
At each stage, all admissible partitions contributing to the same current are
summed coherently before propagation. The final closures join the largest
retained currents to the remaining external states. This ordering makes the
dependency structure explicit without assigning physical meaning to a
particular storage layout.

\subsection{Numerical realization}

Compiled mode combines many relations in a stage into process-specific
multi-output evaluators. Eager mode follows the same process graph using
reusable prepared-model kernels. Recurrence mode represents the current and
closure sequence directly. All three preserve the same coherent sums, colour
contraction, helicity treatment, symmetry weights, and normalization.

For resolved leading-colour evaluation, the public result has the form
\begin{equation}
  R_{p,h,f}=\left|\mathcal M_{h,f}(p)\right|^2,
  \qquad
  \operatorname{shape}(R)=
  (N_{\rm point},N_{\rm helicity},N_{\rm flow}) .
  \label{eq:zgg-resolved-shape}
\end{equation}
Structural zeros remain explicit, and summing the non-point axes must reproduce
the directly evaluated total within the stated numerical tolerance.

\subsection{Timing boundary}

The primary runtime quantity is the complete native wall time after
caller-language array conversion. It includes source and momentum preparation,
the mode-specific numerical evaluation, selectors, closures, and final
reduction. Packing performed by the caller is outside this boundary.

Some backends additionally expose a narrower execution attribution. Because
its precise boundary is mode-specific, it is reported only as a supplementary
quantity and only when measured directly. The complete wall time remains the
common observable used for comparisons; an unavailable attribution is marked
explicitly rather than inferred.
